% =========================================================
% Proof file: prf-unique-additive-identity.tex
% =========================================================

\newpage
\phantomsection
\label{prf:unique-additive-identity}

\begin{remark*}[Return]
\hyperref[thm:unique-additive-identity]{Return to Theorem}
\end{remark*}

\begin{theorem*}[Unique Additive Identity]
If $0,0' \in V$ each satisfy
\[
\forall v \in V,\quad v+0=v
\qquad\text{and}\qquad
\forall v \in V,\quad v+0'=v,
\]
then
\[
0=0'.
\]
\end{theorem*}

\begin{proof}
\textbf{Professional Standard Proof.}~\\
Because $0$ is an additive identity,
\[
0+0'=0'.
\]
Because $0'$ is an additive identity,
\[
0+0'=0.
\]
Hence
\[
0=0'.
\]
Therefore the additive identity is unique.
\end{proof}

\begin{proof}
\textbf{Detailed Learning Proof.}~\\

\textbf{Step 1.} Start from the assumption that both $0$ and $0'$ act as additive identities.\\
This means that for every $v \in V$,
\[
v+0=v
\qquad\text{and}\qquad
v+0'=v.
\]

\textbf{Step 2.} Apply the first identity law to the vector $0'$.\\
Since $0$ is an additive identity,
\[
0'+0=0'.
\]
By commutativity of addition,
\[
0+0'=0'.
\]

\textbf{Step 3.} Apply the second identity law to the vector $0$.\\
Since $0'$ is an additive identity,
\[
0+0'=0.
\]

\textbf{Step 4.} Compare the two values of the same sum.\\
The expression $0+0'$ equals both $0'$ and $0$. Therefore
\[
0=0'.
\]

\textbf{Step 5.} State the conclusion.\\
There cannot be two different additive identities. Hence the additive identity is unique.
\end{proof}

\begin{remark*}[Proof structure]
The argument is direct. The same expression $0+0'$ is evaluated using each identity law, and the two resulting values are equated.
\end{remark*}

\begin{remark*}[Dependencies]
\hfill
\begin{itemize}
  \item \hyperref[def:vector-space]{Definition of Vector Space}
\end{itemize}
\end{remark*}